\section{Prefix-Frontier Entry Discovery}
\label{sec:frontier}
\subsection{Entry Coverage Without Minimal-Superset Reduction}
Let $C_Q=\{S:Q\subseteq S,\ G_S\ne\varnothing\}$ be the eligible observed
label sets.  Write $A\prec_p S$ when the sorted sequence of $A$ is a proper
prefix of that of $S$.  Our entry frontier is
\begin{equation}
 F(Q)=\{S\in C_Q:\nexists A\in C_Q\text{ with }A\prec_p S\}.
\end{equation}
UNG's entry set instead excludes $S$ whenever an eligible proper
\emph{subset} of $S$ is observed.  Every prefix is a subset, but many subsets
are not prefixes.  Consequently the LNG-minimal entry set is contained in
$F(Q)$, and the containment can be strict.  Figure~\ref{fig:running-example}
shows the smallest relevant distinction: $\{b\}\subset\{a,b\}$ is not a
prefix relation under $a<b$, so both groups must remain in $F(\{b\})$.

Connect each terminal to its nearest terminal descendants in the canonical
trie.  This forms a directed forest, with one incoming edge per non-root
terminal.  It has exactly $G-r$ group edges, where $r$ is the number of
terminals without a terminal ancestor.  Thus the count is $G-1$ only if
$r=1$.  The LNG cover graph has at most $O(G^2)$ edges, and two-level subset
families can attain a quadratic count.  This is a worst-case distinction,
not an assertion that every dataset has a dense LNG.

\paragraph{Proposition 1: group-level frontier coverage.}
Every eligible observed terminal is reachable from a terminal in $F(Q)$ by a
path consisting entirely of eligible terminals in the prefix forest.

\emph{Proof.} For an eligible terminal $S$, take the first eligible terminal
$A$ on its root-to-$S$ trie path.  If no earlier eligible terminal exists,
$A\in F(Q)$ by definition.  Consecutive terminals from $A$ to $S$ are joined
by the nearest-terminal edges.  Their label sets all contain $A$, hence
contain $Q$.  This also covers $A=S$ and $Q=\varnothing$. \hfill$\square$

The proposition concerns the label plane.  Lifting it to exhaustive vector
reachability requires retaining every frontier group as a seed, each directed local graph being reachable from its seeded or incoming
portals, and accessible materialized cross edges
for each required group transition.  Finite candidate capacity, approximate
local graphs, pruning, and seed selection can violate those additional
conditions.  We therefore measure Recall; the proposition is not an ANN
accuracy guarantee.

\subsection{The Measured Entry Algorithm}
The measured Trie provider stores, for each label $a$, a compressed bitmap
$I_a$ of groups containing $a$, and for each terminal $S$, its proper terminal
ancestors $\mathrm{Anc}_p(S)$.  It intersects postings to obtain $C_Q$ and
tests only prefix ancestors to produce $F(Q)$
(Algorithm~\ref{alg:frontier}).  This algorithm avoids comparing $S$ against
eligible terminals on unrelated branches.  The bitmaps and ancestor lists
are initialized before steady-state timing; per-query intersections and
ancestor tests remain in the discovery time.

\begin{algorithm}[t]
\caption{Bitmap-backed prefix frontier (measured provider)}
\label{alg:frontier}
\begin{algorithmic}[1]
\Require query labels $Q$, label bitmaps $I_a$, proper terminal ancestors
\If{$Q=\varnothing$}
  \State \Return precomputed terminal roots of the prefix forest
\EndIf
\If{any query label has no posting}
  \State \Return $\varnothing$
\EndIf
\State copy the smallest query-label posting into $C_Q$
\For{each remaining query label $a$}
  \State $C_Q\gets C_Q\cap I_a$
\EndFor
\State $F\gets\varnothing$
\For{each $S\in C_Q$}
  \If{no $A\in\mathrm{Anc}_p(S)$ belongs to $C_Q$}
    \State append the group identifier of $S$ to $F$
  \EndIf
\EndFor
\State \Return $F$
\end{algorithmic}
\end{algorithm}

For comparison, both LNG-semantic providers first use a trie walk to collect
a prefix-minimal candidate frontier, rather than enumerate all eligible
supersets. They then reduce this frontier under set containment.  The original version sorts candidates by label cardinality and
removes sets containing an already retained minimum.  The optimized version
uses cardinality buckets when the size span is small, with a sorting fallback;
it still performs pairwise sorted-set containment checks for minimality.
It preserves the original entry definition.  The prefix provider changes
that definition and the data structures used to compute it.  Thus its timing
advantage is not solely the effect of removing one sort.

\paragraph{Cost and output tradeoff.}
Let $c=|C_Q|$, $f=|F(Q)|$, and $A_Q$ be the number of proper-ancestor
membership probes actually performed.  If $T_{\cap}$ is the cost of copying
and intersecting compressed postings and $t_{\in}$ bounds one bitmap
membership probe, prefix discovery costs
\begin{equation}
 O\bigl(|Q|+T_{\cap}+c+A_Qt_{\in}+f\bigr).
\end{equation}
For maximum label length $h$, $A_Q\le ch$.  Bitmap representation and
compression affect $T_{\cap}$ and $t_{\in}$; treating both as unit-time set
operations would hide substantial work.  The stored ancestor lists require
up to $O(Gh)$ entries, in addition to postings and trie storage.  The empty
predicate uses the precomputed root frontier in $O(r)$ output time.

For the LNG providers, let $c_{\mathrm{LNG}}$ be the number of frontier candidates,
$T_{\mathrm{cand}}$ their discovery time, and $P_Q$ the number of pairwise minimality tests.  With sorted label sets of
length at most $h$, their cost is bounded by
\begin{equation}
 O\bigl(T_{\mathrm{cand}}+\mathrm{Order}(c_{\mathrm{LNG}})+P_Qh\bigr),
 \qquad P_Q\le c_{\mathrm{LNG}}(c_{\mathrm{LNG}}-1)/2,
\end{equation}
where ordering is $O(c_{\mathrm{LNG}}\log c_{\mathrm{LNG}})$ for sorting or
$O(c_{\mathrm{LNG}}+\Delta)$ for cardinality
buckets over span $\Delta$.  The
prefix algorithm restricts which relations are tested, but may return more
groups and hence increase seed distance work and graph traversal.  We report
these downstream costs separately.
